\documentclass{article}

\title{Ch. 24 HW \#2 Solutions}
\author{Principles of Macroeconomics}
\date{\today}

\begin{document}

\maketitle

\section*{Questions for Review}

\begin{enumerate}

\setcounter{enumi}{4}
\item Consumption: A household buys a new blender. \\
	Investment: A company buys a new forklift. \\
	Government Purchases: The government buys a new tank. \\
	Net Exports: A person in Great Britain buys Microsoft Office
	
\item Unlike nominal GDP, real GDP isn't influenced by changes in the price level. Therefore if you compare real GDP across periods, it identifies the change in an economy's productive capacity with no distortions from changing prices. Productive capacity is a better measure of an economy's well-being because it determines the economy's ability to meet people's material needs.

\item \phantom{x}
	
	\begin{center}
	\begin{tabular}{| c | c | c | c |}
	\hline
	Year & $GDP_{Nom}$ & $GDP_{Real}$ & GDP Deflator \\
	\hline
	2023 & $100(2) = 200$ & $100(2) = 200$ & $\frac{200}{200} \times 100 = 100$ \\
	\hline
	2024 & $200(3) = 600$ & $200(2) = 400$ & $\frac{600}{400} \times 100 = 150$ \\
	\hline
	\end{tabular}
	\end{center}

\end{enumerate}

\section*{Problems \& Applications}

\begin{enumerate}

\item

	\begin{enumerate}
	
	\item Consumption -- A personal refrigerator is household spending.
	
	\item Investment -- New housing is investment.
	
	\item None -- There's no spending on new production.
	
	\item Consumption -- A personal haircut is household spending.
	
	\item Consumption and investment -- Purchasing a personal car increases consumption, but decreasing inventory decreases investment by an offsetting amount.
	
	\item Investment -- A new car is a capital good for a rental car company because it allows them to produce more rental agreements.
	
	\item Government purchases -- Repaving roads requires government spending on current production.
	
	\item None -- Transfer payments aren't spending on production. 
	
	\item Consumption and net exports -- Buying a foreign bottle of wine increases consumption, but buying an import decreases net exports by an offsetting amount.
	
	\item Investment -- A factory is a capital good. 
	
	\end{enumerate}
	
\item \phantom{x}

	\begin{center}
	\begin{tabular}{| c | c | c | c |}
	\hline
	Year & $GDP_{Real}$ & $GDP_{Nom}$ & GDP Deflator \\
	\hline
	1970 & 3,000 & 1,200 & $\frac{1,200}{3,000}(100) = 40$ \\
	\hline
	1980 & 5,000 & $5,000 \left( \frac{60}{100} \right) = 3,000$ & 60 \\
	\hline
	1990 & $\frac{6,000}{100} \times 100 = 6,000$ & 6,000 & 100 \\
	\hline
	2000 & 8,000 & 8,000 & 100 \\
	\hline
	2010 & $\frac{15,000}{200} \times 100 = 7,500$ & 15,000 & 200 \\
	\hline
	2020 & 10,000 & $\left( \frac{300}{100} \right) = 30,000$ & 300 \\
	\hline
	2030 & 20,000 & 50,000 & $\frac{50,000}{20,000}(100) = 250$ \\
	\hline
	\end{tabular}
	\end{center}
	
\item GDP is spending on current production. Transfer payments aren't made in exchange for production, so they are excluded.

\end{enumerate}

\end{document}